\chapter{اولویت‌های تقنینی فاز اول}

مجلس مهستان در فاز اول (صد روز نخست) موظف است اقدامات تقنینی زیر را به ترتیب اولویت انجام دهد.

\section{فرمان‌های اضطراری (هفته‌ی اول)}
این اقدامات بلافاصله و بدون نیاز به فرایند عادی صادر خواهند شد.

\begin{tabular}{clp{8cm}}
\toprule
\textbf{ردیف} & \textbf{اقدام} & \textbf{دلیل فوریت} \\
\midrule
۱ & آزادی زندانیان سیاسی & حق بنیادین؛ نمادین \\
۲ & لغو حجاب اجباری & رفع تبعیض و نماد سرکوب \\
۳ & رفع فیلترینگ اینترنت & حق اطلاعات و زیرساخت دموکراسی \\
۴ & تعلیق مجازات اعدام & حق حیات و اجماع بین‌المللی \\
۵ & آزادی فعالیت احزاب & حق تجمع و تشکل \\
۶ & حفاظت از اسناد امنیتی & مستندسازی عدالت انتقالی \\
\bottomrule
\end{tabular}

\section{مصوبات اولویت بالا (هفته‌ی ۱ تا ۴)}
\begin{itemize}
    \item \textbf{قانون آزادی مطبوعات:} تضمین آزادی بیان نهادینه.
    \item \textbf{قانون اجتماعات:} آزادی تظاهرات بدون نیاز به مجوز پیشینی.
    \item \textbf{قانون شفافیت مالی:} الزام مقامات به انتشار دارایی‌ها.
\end{itemize}

\section{ترکیب تخصص و دموکراسی مستقیم}
فرایند تصویب هر قانون از الگوی دو بازویی پیروی می‌کند:
\begin{enumerate}
    \item \textbf{مرکز تحقیقات (بازوی تکنوکراتیک):} بررسی داده‌ها، آمار و تجربه‌ی جهانی.
    \item \textbf{واحد ارتباط مردمی (بازوی دموکراتیک):} نظرسنجی، بازخورد و رفراندوم موضوعی.
\end{enumerate}

\begin{center}
\begin{tikzpicture}[node distance=1.5cm]
    \node (unit) [draw, fill=LightBlue, text width=4cm, align=center, rounded corners] {بازوی دموکراتیک \\ (صدای مردم)};
    \node (research) [draw, fill=SoftPurple!30, text width=4cm, align=center, rounded corners, right=2cm of unit] {بازوی تکنوکراتیک \\ (واقعیت و تخصص)};
    \node (assembly) [draw, fill=SoftGreen, text width=6cm, align=center, rounded corners, below=2cm of $(unit.south)!0.5!(research.south)$] {مجلس مهستان \\ (تصمیم‌گیری آگاهانه)};
    
    \draw[->, thick] (unit) -- (assembly);
    \draw[->, thick] (research) -- (assembly);
\end{tikzpicture}
\end{center}
